\section{Exact binary sumsets}

Which integers are sums of two primitive values of a quadratic norm?
We determine the complete sumsets in four class-number-one cases,
including every small exception and a proof covering every remaining
target.

The companion paper \cite{pointwisecompanion} develops broader pointwise
correlation asymptotics. The exact classifications here have their own
classical proof route and do not depend on that paper.

For a negative discriminant $d$, put
\[
 P_d=\{Q_d(x,y)>0:x,y\in\Z,\ \gcd(x,y)=1\},\qquad
 \begin{array}{c|c}
 d&Q_d(x,y)\\\hline
 -3&x^2-xy+y^2\\
 -7&x^2+xy+2y^2\\
 -8&x^2+2y^2\\
 -11&x^2+xy+3y^2.
 \end{array}
\]
Zero coordinates and both signs are allowed.  Each summand is a positive
numerical value; different coordinate representations do not create
new summands.

\begin{theorem}\label{thm:exact-overview}
The exact binary sumsets are
\[
\begin{array}{c|l}
 d&P_d+P_d\\\hline
 -3&\{N\ge2:N\text{ even},\ N\not\equiv0,12\pmod{18}\},\\
 -7&\{N\ge2:49\nmid N\}\setminus\{7,77,203\},\\
 -8&\{N\ge2:16\nmid N\},\\
 -11&\{N\ge2:N\text{ even},\ 121\nmid N\}.
\end{array}
\]
\end{theorem}

For the Eisenstein row, let
\[
 \AA=\{n\ge1:p\mid n\Longrightarrow p\equiv1\pmod3\}.
\]
Here \(1\in\AA\). Proposition~\ref{prop:primitive-classification} proves
\[
                    P_{-3}=\AA\sqcup(3\!\cdot\!\AA).
\]
An inert prime cannot divide a primitive norm, while the ramified
prime occurs at most once; conversely the split-prime core and its
ramified dilation have primitive witnesses. The same reduction in
the other three rings explains why we prove stronger core sumset
identities before restoring ramified factors.

Our method adapts the quaternionic norm-shell and torus-normalizer
argument used for the Gaussian restricted-prime theorem of Howe,
Huffman, Kohlmeyer and Kruer \cite[Theorem~1.1]{Euler}.  A quaternion
shell produces many candidate pairs of binary norms.  The Hecke
spectrum bounds the fraction spoiled by each forbidden prime.
Exact dyadic and ramified filters keep the total fraction below one.
The spectral input is classical Jacquet--Langlands and the weight-two
Ramanujan bound; Section~\ref{sec:exact-hecke} states its normalization.
The applications require different local quotients, including all ideal
classes and inverse-unit mass at discriminant eleven.

Explicit estimates settle an infinite range; rational interval
certificates meet complete finite scans, and multiplicative descent
exhausts the targets. Appendix~\ref{app:verifier} explains the exact
programs and completeness arguments.
Table~\ref{tab:exact-coverage} records these joins. Here \(C_d\) is
the core defined in its section, \(n\) is a reduced target, and
\(m\) at \(-8\) is an odd inert-supported core. Each interval
certificate covers every target of the indicated reduced families in its interval.

\begin{table}[ht]
\centering\small
\setlength{\tabcolsep}{3pt}
\renewcommand{\arraystretch}{1.13}
\caption{Coverage of the exact classifications.  All finite scans include
only the locally eligible targets of the indicated families.}
\label{tab:exact-coverage}
\begin{tabular}{@{}>{\raggedright\arraybackslash}p{.045\linewidth}>{\raggedright\arraybackslash}p{.17\linewidth}>{\raggedright\arraybackslash}p{.255\linewidth}>{\raggedright\arraybackslash}p{.225\linewidth}>{\raggedright\arraybackslash}p{.24\linewidth}@{}}
\toprule
$d$&Complete finite scan&Rational interval bridge&Infinite estimate&Descent and dilation\\
\midrule
$-3$&$N<10^8$: $\AA+\AA$ and $\AA+3\!\cdot\!\AA$.
&No interval bridge needed.
&Reduced $n\ge10^8$, in both admissible classes modulo six.
&Remove all split-prime factors, then scale back; use $P_{-3}=\AA\sqcup3\!\cdot\!\AA$.\\[3pt]
$-7$&$N\le2^{24}$: $C_7+C_7$ and $P_7+P_7$, with their listed exceptions.
&$257$ intervals on $[2^{24},2^{37})$; $17$ on $[2^{37},2^{54})$, for inert-supported targets and split primes.
&The same reduced families for $n\ge2^{54}$.
&Remove split factors; patch quotients $11,29$.  Ramified dilation by seven gives the primitive exceptions.\\[3pt]
$-8$&$N\le2^{27}$: $C_8+C_8$ and $C_8+2C_8$.
&$107$ intervals: $m\in[2^{27},2^{60})$ for mixed sums, $[2^{26},2^{60})$ at $2m$, and $[2^{20},2^{60})$ at $4m$.
&Inert cores $m\ge2^{60}$; the mixed prime boundary holds for primes $N\ge2^{22}$.
&Remove split factors; the remaining targets are $m,2m,4m$.  Use $P_{-8}=C_8\sqcup2C_8$.\\[3pt]
$-11$&$N\le2^{25}$: $C_{11}+C_{11}$.
&$35$ intervals on $[2^{25},2^{60})$, for even inert-supported targets prime to eleven.
&The same reduced family for $n\ge2^{60}$.
&Remove split factors; ramified dilation by eleven gives $P_{-11}+P_{-11}$.\\
\bottomrule
\end{tabular}
\end{table}

Section~\ref{sec:eisenstein} proves the two Eisenstein core identities;
Sections~\ref{sec:minus7}--\ref{sec:minus11} treat the other discriminants.
